\documentclass[11pt,a4paper]{article}
\usepackage[T1]{fontenc}
\usepackage{lmodern,amsmath,amssymb,amsthm,mathtools,booktabs,microtype}
\usepackage[margin=27mm]{geometry}
\usepackage[colorlinks=true,linkcolor=blue,citecolor=blue,urlcolor=blue]{hyperref}
\setlength{\emergencystretch}{2em}
\newtheorem{theorem}{Theorem}[section]
\newtheorem{lemma}[theorem]{Lemma}
\newtheorem{corollary}[theorem]{Corollary}
\newtheorem{proposition}[theorem]{Proposition}
\newcommand{\Q}{\mathbb Q}\newcommand{\Z}{\mathbb Z}
\title{Finite square-class conditions for native detector inputs}
\author{}\date{16 November 2026}
\begin{document}\maketitle
\begin{abstract}
We give a finite-square-class version of the rational-complement
criterion for the general and conditional elliptic detector sources.
Finitely many complement levels and square classes supported on the
original exceptional primes require no additional retained coordinates.
All complement norms are automatically nonzero at prepared actual
point tuples. A corresponding rational realization theorem would yield
twenty-one or nine total variables; that realization remains an explicit
unproved premise. An infinite family of prepared coefficient inputs
demonstrates, for exceptional sets allowing the class $409$, that the
square-class condition is strictly weaker than rational squareness,
without asserting realization by native point tuples.
\end{abstract}

\section{The fixed exceptional set}
Fix the general source or either conditional source covered
by~\cite{Complements}. Retain its arithmetic premises, a coefficient
member in its sufficient soundness range, and its original fixed set
$S$ of exceptional rational primes. In particular $17\in S$.
Every prime at which the source requires its negative-odd-valuation
implication remains outside $S$. The set is not enlarged during the
construction; the fixed consumed prime of the general source is not
made exceptional.

Let $\mathcal M\subset\Z_{>0}$ and $\mathcal D\subset\Q_{>0}$ be
nonempty finite sets, with
\begin{equation}\label{eq:classes}
 v_p(d)\in2\Z\qquad(d\in\mathcal D,\ p\notin S).
\end{equation}
We may replace each $d$ by the positive squarefree integer representing
its rational square class. Thus $\mathcal D$ may be taken to consist of
distinct products of primes in $S$. Put
\[
 a_m=17m^2,\qquad C_m(b)=a_m(a_m-b)\quad(m\in\mathcal M).
\]
Both finite sets are fixed coefficient data, not quantified variables.

\begin{lemma}\label{lem:local}
If $dC_m(b)$ is a rational square for some $m\in\mathcal M$ and
$d\in\mathcal D$, then $b$ has no negative odd valuation at any
prime outside $S$.
\end{lemma}
\begin{proof}
At such a prime, $v_p(a_m)=2v_p(m)$ is nonnegative and even. If
$v_p(b)<0$ were odd, it would be strictly smaller than $v_p(a_m)$.
Consequently
\[
 v_p\bigl(d a_m(a_m-b)\bigr)
       =v_p(d)+v_p(a_m)+v_p(b)
\]
would be odd, contradicting squareness. If $C_m(b)=0$, then $b=a_m$
has no negative valuation and satisfies the conclusion directly.
\end{proof}

\section{Guards at every point-algebra character}
At an actual prepared point-stage tuple, the finite point algebra
splits over $K=\Q(\sqrt2)$. Each detector input has the form
\begin{equation}\label{eq:native}
 b=wT^2,
\end{equation}
where $w\in\Q^\times$ is independent of the character and $T$ is a
unit of the point algebra. At the selected rational character, $T$ is
a nonzero rational number. Preparation makes $v_{17}(b)$ even, hence
$v_{17}(w)$ is even. These facts, including all characters rather than
only the selected one, are proved in~\cite[Sections 1--2]{Complements}.

\begin{lemma}\label{lem:guard}
At every such tuple, $a_m-b$ is a unit of the whole point algebra
for every $m\in\mathcal M$.
\end{lemma}
\begin{proof}
If a geometric character $\chi$ annihilated $a_m-b$, then
\[
                    a_m/w=\chi(T)^2\in K^{\times2}.
\]
A nonzero rational number that is a square in $K$ belongs to
$\Q^{\times2}$ or $2\Q^{\times2}$: the irrational coefficient of
$(u+v\sqrt2)^2$ vanishes only when $uv=0$. Both classes have even
valuation at $17$, whereas
\[
             v_{17}(a_m/w)=1+2v_{17}(m)-v_{17}(w)
\]
is odd. This is impossible. An element of a finite algebra over a
field is a unit exactly when it is nonzero at every geometric character.
\end{proof}

The conclusion holds simultaneously for the entire fixed finite set.
For a prepared input $0<b<17$, every $a_m\ge17$, so every $C_m(b)$
is positive as well as nonzero.

\section{The conditional variable reduction}
The realization condition below concerns the unchanged native detector
functions, not freely chosen coefficient inputs.

\begin{theorem}\label{thm:reduction}
Suppose that for every integer input required by the chosen source's
completeness statement, there is an actual rational point-stage tuple
with all inherited guards nonzero, prepared inputs $b_0,b_1$, and choices
$m_i\in\mathcal M$, $d_i\in\mathcal D$ such that
\begin{equation}\label{eq:realization}
                         d_iC_{m_i}(b_i)\in\Q^2
                         \qquad(i=0,1).
\end{equation}
Then, under the retained source premises, there is an integer polynomial
defining $\Z$ over $\Q$ in twenty-one total variables for the general
source, or nine total variables for either conditional source. The input
is included. Condition~\eqref{eq:realization} is an additional unproved
arithmetic premise, not a conclusion of the guard lemmas.
\end{theorem}
\begin{proof}
First fix all four finite choices. Write $b_i=N_i/D_i$ in the native
point algebra, with each $D_i$ a unit by the inherited guards. Adjoin
the two finite roots
\begin{equation}\label{eq:roots}
 Q_i^2=d_i a_{m_i}D_i(a_{m_i}D_i-N_i),\qquad i=0,1.
\end{equation}
For squarefree integer representatives $d_i$, these are monic equations
over the integral point algebra. At a rational character, $Q_i/D_i$ is
a square root of $d_iC_{m_i}(b_i)$, and conversely every such square
root extends that character. Include the norms of
$a_{m_i}D_i-N_i$ in the scalar guard, along with all inherited point
and denominator guards.

On this guarded locus the two radicands are units. The finite algebra
is therefore etale whenever the inherited point algebra is etale, as
required by the separating-element construction of~\cite{UnitNorm}.
It has rank $4^{p+1}$ when there are $p$ point pairs. The polynomial
denominator envelopes, norm conditions and single separating coordinate
are exactly those of the guarded construction in~\cite{Complements}.
Lemma~\ref{lem:local} gives the same local implication used by source
soundness, at every required prime. All other soundness arguments are
unchanged. Lemma~\ref{lem:guard} makes the added norms automatically
nonzero at the actual tuple supplied by~\eqref{eq:realization}.

The retained coordinates are the input, $2p$ point coordinates and
the comparison slope. The encoder adds one, giving $2p+3$: this is
twenty-one for $p=9$ and nine for $p=3$.

Finally construct one such polynomial for each pair of finite choices
and take their product, using the same variable names, including the
same encoder variable. Over $\Q$ the product vanishes exactly when
one factor vanishes. This finite union does not increase the variable
count. For a conditional source, the usual input factor includes zero
and need occur only once after the finite product is formed.
\end{proof}

The product construction does not require distinct radicands from
different choices. Coincident radicands can occur, and treating all
choices as an unlabelled quadratic extension would not justify a
distinct-character separator. Separate encodings avoid that issue.
If there are $r=|\mathcal M||\mathcal D|$ choices per block and every
individual polynomial has degree at most $\Delta$, the product has
degree at most $r^2\Delta$, with the single input factor accounted for
separately when applicable. No numerical degree or coefficient bound
is assigned here.

\begin{corollary}\label{cor:allclasses}
Take $\mathcal M=\{1\}$ and let $\mathcal D$ contain all positive
squarefree products of primes in $S$. Then the realization premise
in Theorem~\ref{thm:reduction} is equivalent to requiring both actual
complements $17(17-b_i)$ to have even valuation at every prime outside
$S$. There are $2^{|S|}$ possible positive square classes.
\end{corollary}
\begin{proof}
Each complement is positive. A positive rational number whose valuations
outside $S$ are all even is a rational square times a unique squarefree
product $d$ of primes in $S$. Multiplying by that same $d$ makes it a
rational square. The converse follows immediately from
\eqref{eq:classes}.
\end{proof}

\section{Prepared inputs with a nonsquare complement}
The following family establishes strict enlargement on the coefficient
locus. It is not asserted to arise from actual native detector tuples.

\begin{proposition}\label{prop:family}
For $k\in\Z_{\ge0}$ put $t=10+187272k$ and
\[
 s(t)=\frac{34(3272t^2-409t-136)}{69(17+409t^2)},\qquad
 q(t)=\frac{17(17-544t-409t^2)}{69(17+409t^2)},
 \qquad b(t)=s(t)^2.
\]
These give infinitely many distinct prepared positive rational inputs
with
\[
 0<b(t)<17,\qquad
 \bigl(v_2(b),v_3(b),v_{17}(b)\bigr)=(4,2,2),
 \qquad17(17-b(t))=409q(t)^2\ne0.
\]
Their reduced denominators are squares, and the inputs are squares in
$\Q_2$, $\Q_3$ and $\Q_{17}$. Their complements are not rational
squares. If $409\in S$, their complement classes are permitted in
Corollary~\ref{cor:allclasses}.
\end{proposition}
\begin{proof}
The displayed functions parametrize the conic
$17s^2+409q^2=289$ from the point
$(s,q)=(272/69,17/69)$. Direct substitution proves the equation.
For $t\ge10$, the numerator of $s(t)$ is positive and
\[
 \frac{272}{69}-s(t)
   =\frac{2(17\cdot272+409\cdot17t)}{69(17+409t^2)}>0.
\]
Since $(272/69)^2<17$, the required real inequalities follow.
The numerator of $q(t)$ is negative for $t\ge10$, so it is nonzero.

At $t=10$ the values are
\[
 s=\frac{1220124}{313697},\qquad
 q=-\frac{87499}{313697},\qquad
 b=\frac{1488702575376}{98405807809}.
\]
In the unreduced numerator and denominator polynomials of $s$, their
valuation pairs at $2,3,17$ are respectively $(2,0)$, $(3,2)$ and
$(1,0)$. Now $187272=2^3 3^4 17^2$. Substituting
$t=10+187272k$ leaves each polynomial congruent to its value at $10$
modulo a strictly higher power of the respective prime. The valuations
of $s$ are therefore always $(2,1,1)$, giving the stated values for $b$.
All other preparation assertions follow from $b=s^2$.

The rational function $s(t)^2$ is nonconstant, so its values on this
infinite progression are infinite in number. The integer $409$ is prime
and is not a rational square; since $q(t)\ne0$, no displayed complement
is a rational square.
\end{proof}

The missing step for a defining-variable reduction remains actual
rational realization in Theorem~\ref{thm:reduction}. A family of prepared
coefficient inputs, even one with these exact class constraints, does
not supply equality with the native reciprocal traces.

\begin{thebibliography}{9}
\bibitem{Complements}
\emph{Unit guards and a finite square-complement reduction for elliptic detectors},
project source,
\nolinkurl{papers/pole_parameters/guarded_complement_reduction_20261116_v45.tex}.
\bibitem{UnitNorm}
\emph{A single norm determinant for the 25- and 13-variable constructions},
project source,
\nolinkurl{papers/finite_algebra/unit_norm_20261114_v28.tex}.
\end{thebibliography}
\end{document}
